\section{Benchmark Results}
\label{sec:results}
\subsection{Existing Evidence}
The existing source-fixed port was checked on two nine-scene, 60-second subsets against stored scores. These checks establish numerical reproduction on those inputs, not construct validity. Existing 60-second image-conditioned scores permit matched comparison of source-fixed SNF, VBench-I2V and VBench-Long on shared source scenes. The 120- and 240-second method records score 60-second endpoint windows of separately requested generations; these should not be interpreted as trajectories of a single rollout. Historical text-to-video records use different mask and measurement identities and are reported separately.

\begin{table*}[t]
\caption{\textbf{Coverage of existing development evidence.} Counts are scenes or prompts per method, not independent confirmation samples. Text- and image-conditioned protocols use different masks and references; their numbers should not be pooled. Some samples are also used in simultaneous method submissions and are not used as new qualitative examples here.}
\label{tab:existing_evidence}
\centering
\setlength{\tabcolsep}{4pt}
\begin{tabular}{@{}lllll@{}}
\toprule
Track & 60 s & 120 s & 240 s & Interpretation \\
\midrule
Text to video & 22 shared prompt IDs & 6 prompts & 4 prompts & Legacy masks; prompt sets differ across durations \\
Image to video & 48 appearance / 41 motion & 48 / 41 & 48 / 41 & Shared source masks; 60 s endpoint windows \\
\bottomrule
\end{tabular}
\end{table*}

Table~\ref{tab:existing_evidence} describes the available development records rather than a new leaderboard. The text-conditioned 60 s count is the unambiguous shared-identifier intersection in the existing comparison; identifiers still need video-hash verification. The 120 and 240 s counts are for the seven historical checkpoints. The image-conditioned counts describe paired examples from the previously scored family; different measures can have different eligible populations. The image-conditioned endpoint comparisons use separately requested durations, not prefixes of one long video.

\begin{figure*}[t]
\centering
\setlength{\tabcolsep}{1.1pt}
\newcommand{\iifigpath}{figures/i2v_comparison_2026_11_08}
\newcommand{\iicell}[2]{\includegraphics[width=.191\textwidth]{\iifigpath/#1_#2_end.png}}
\begin{tabular}{@{}ccccc@{}}
\scriptsize Source & \scriptsize Causal-Forcing & \scriptsize LTX-Video & \scriptsize Wan2.1 & \scriptsize Wan2.2 \\
\includegraphics[width=.191\textwidth]{\iifigpath/vb161950_source.png} &
\iicell{vb161950}{causal_forcing} & \iicell{vb161950}{ltx_video} &
\iicell{vb161950}{wan21} & \iicell{vb161950}{wan22} \\
\includegraphics[width=.191\textwidth]{\iifigpath/vb10462563_source.png} &
\iicell{vb10462563}{causal_forcing} & \iicell{vb10462563}{ltx_video} &
\iicell{vb10462563}{wan21} & \iicell{vb10462563}{wan22} \\
\end{tabular}
\caption{\textbf{Image-conditioned comparison on two sources.} Source images and the final available frame of one approximately 5\,s output per system. The waterfall and campfire illustrate different moving materials; each row shares the same input image and prompt across systems. The four runs use their documented inference pipelines, with configuration differences recorded in the accompanying provenance. These short examples do not establish long-horizon performance, motion realism or a ranking. Sources: VBench++ image suite~\cite{huang2024vbench++}, using Pexels images \href{https://www.pexels.com/photo/161950/}{161950} and \href{https://www.pexels.com/photo/10462563/}{10462563}; images are cropped and resized.}
\label{fig:i2v_comparison}
\end{figure*}


Figure~\ref{fig:i2v_comparison} illustrates image-conditioned outputs from four released pipelines on two sources. These are five-second examples, not the planned minute-scale comparisons. The prior 240-second text-conditioned figure reused rollout samples shown in a simultaneous Steady-Forcing submission and has been removed. \pending{non-overlapping text-to-video qualitative comparison and longer image-to-video comparisons}

\subsection{Leaderboard}
Results are reported per track and horizon as per-axis profiles with 95\% source-clustered bootstrap intervals~\cite{efron1979bootstrap}, macro-averaged over regimes, with eligibility counts beside every value. No single summed score is reported. \pending{leaderboard tables per track and horizon}

\subsection{Per-Regime Analysis and Comparison with Established Metrics}
\pending{per-regime breakdown; agreement with VBench and WorldScore dimensions}

\pending{rights-cleared independent confirmation set, independently reviewed role maps, native long-video baselines, sufficient multi-seed coverage, real-periodicity controls and human construct validation. Previously scored videos and metrics are reused rather than regenerated or rescored where their input and protocol identity is verified.}
